\begin{thebibliography}{99}

\bibitem[Bayne(2010)]{bayne2010}
 Bayne, T. (2010). \textit{The Unity of Consciousness}. Oxford University Press. \href{https://doi.org/10.1093/acprof:oso/9780199215386.001.0001}{doi:10.1093/acprof:oso/9780199215386.001.0001}.

\bibitem[Caminiti et al.(2009)]{caminiti2009}
 Caminiti, R., Ghaziri, H., Galuske, R., Hof, P. R., \& Innocenti, G. M. (2009). Evolution amplified processing with temporally dispersed slow neuronal connectivity in primates. \textit{Proceedings of the National Academy of Sciences}, 106(46), 19551--19556. \href{https://doi.org/10.1073/pnas.0907655106}{doi:10.1073/pnas.0907655106}.

\bibitem[Caminiti et al.(2013)]{caminiti2013}
 Caminiti, R., Carducci, F., Piervincenzi, C., Battaglia-Mayer, A., Confalone, G., Visco-Comandini, F., Pantano, P., \& Innocenti, G. M. (2013). Diameter, length, speed, and conduction delay of callosal axons in macaque monkeys and humans: Comparing data from histology and magnetic resonance imaging diffusion tractography. \textit{Journal of Neuroscience}, 33(36), 14501--14511. \href{https://doi.org/10.1523/JNEUROSCI.0761-13.2013}{doi:10.1523/JNEUROSCI.0761-13.2013}.

\bibitem[Chung \& Fuchs(1951)]{chung1951}
 Chung, K. L., \& Fuchs, W. H. J. (1951). On the distribution of values of sums of random variables. \textit{Memoirs of the American Mathematical Society}, 6, 1--12.

\bibitem[de Haan et al.(2020)]{dehaan2020}
 de Haan, E. H. F., Corballis, P. M., Hillyard, S. A., Marzi, C. A., Seth, A., Lamme, V. A. F., Volz, L., Fabri, M., Schechter, E., Bayne, T., Corballis, M. C., \& Pinto, Y. (2020). Split-brain: What we know now and why this is important for understanding consciousness. \textit{Neuropsychology Review}, 30(2), 224--233. \href{https://doi.org/10.1007/s11065-020-09439-3}{doi:10.1007/s11065-020-09439-3}.

\bibitem[Doyle \& Snell(1984)]{doyle1984}
 Doyle, P. G., \& Snell, J. L. (1984). \textit{Random Walks and Electric Networks}. Carus Mathematical Monographs 22. Mathematical Association of America.

\bibitem[Findlay et al.(2024)]{findlay2024}
 Findlay, G., Marshall, W., Albantakis, L., David, I., Mayner, W. G. P., Koch, C., \& Tononi, G. (2024). Dissociating artificial intelligence from artificial consciousness. \textit{arXiv:2412.04571}. \href{https://doi.org/10.48550/arXiv.2412.04571}{doi:10.48550/arXiv.2412.04571}.

\bibitem[Francis et al.(2003)]{francis2003}
 Francis, J. T., Gluckman, B. J., \& Schiff, S. J. (2003). Sensitivity of neurons to weak electric fields. \textit{Journal of Neuroscience}, 23(19), 7255--7261. \href{https://doi.org/10.1523/JNEUROSCI.23-19-07255.2003}{doi:10.1523/JNEUROSCI.23-19-07255.2003}.

\bibitem[Fr\"{o}hlich \& Spencer(1981)]{frohlich1981}
 Fr\"{o}hlich, J., \& Spencer, T. (1981). The Kosterlitz--Thouless transition in two-dimensional Abelian spin systems and the Coulomb gas. \textit{Communications in Mathematical Physics}, 81(4), 527--602. \href{https://doi.org/10.1007/BF01208273}{doi:10.1007/BF01208273}.

\bibitem[Fr\"{o}hlich \& McCormick(2010)]{frohlich2010}
 Fr\"{o}hlich, F., \& McCormick, D. A. (2010). Endogenous electric fields may guide neocortical network activity. \textit{Neuron}, 67(1), 129--143.

\bibitem[Gratiy et al.(2017)]{gratiy2017}
 Gratiy, S. L., Halnes, G., Denman, D., Hawrylycz, M. J., Koch, C., Einevoll, G. T., \& Anastassiou, C. A. (2017). From Maxwell's equations to the theory of current-source density analysis. \textit{European Journal of Neuroscience}, 45(8), 1013--1023. \href{https://doi.org/10.1111/ejn.13534}{doi:10.1111/ejn.13534}.

\bibitem[Herzog et al.(2016)]{herzog2016}
 Herzog, M. H., Kammer, T., \& Scharnowski, F. (2016). Time slices: What is the duration of a percept? \textit{PLoS Biology}, 14(4), e1002433. \href{https://doi.org/10.1371/journal.pbio.1002433}{doi:10.1371/journal.pbio.1002433}.

\bibitem[Hinton(2022)]{hinton2022}
 Hinton, G. (2022). The forward-forward algorithm: Some preliminary investigations. \textit{arXiv:2212.13345}. \href{https://doi.org/10.48550/arXiv.2212.13345}{doi:10.48550/arXiv.2212.13345}.

\bibitem[Hong et al.(2005)]{hong2005}
 Hong, H., Park, H., \& Choi, M. Y. (2005). Collective synchronization in spatially extended systems of coupled oscillators with random frequencies. \textit{Physical Review E}, 72(3), 036217. \href{https://doi.org/10.1103/PhysRevE.72.036217}{doi:10.1103/PhysRevE.72.036217}.

\bibitem[Jacobson et al.(2005)]{jacobson2005}
 Jacobson, G. A., Diba, K., Yaron-Jakoubovitch, A., Oz, Y., Koch, C., Segev, I., \& Yarom, Y. (2005). Subthreshold voltage noise of rat neocortical pyramidal neurones. \textit{The Journal of Physiology}, 564(1), 145--160. \href{https://doi.org/10.1113/jphysiol.2004.080903}{doi:10.1113/jphysiol.2004.080903}.

\bibitem[Johnson(1928)]{johnson1928}
 Johnson, J. B. (1928). Thermal agitation of electricity in conductors. \textit{Physical Review}, 32(1), 97--109. \href{https://doi.org/10.1103/PhysRev.32.97}{doi:10.1103/PhysRev.32.97}.

\bibitem[Kajikawa \& Schroeder(2011)]{kajikawa2011}
 Kajikawa, Y., \& Schroeder, C. E. (2011). How local is the local field potential? \textit{Neuron}, 72(5), 847--858. \href{https://doi.org/10.1016/j.neuron.2011.09.029}{doi:10.1016/j.neuron.2011.09.029}.

\bibitem[Katz \& Miledi(1965)]{katz1965}
 Katz, B., \& Miledi, R. (1965). The measurement of synaptic delay, and the time course of acetylcholine release at the neuromuscular junction. \textit{Proceedings of the Royal Society B}, 161(985), 483--495. \href{https://doi.org/10.1098/rspb.1965.0016}{doi:10.1098/rspb.1965.0016}.

\bibitem[Keller et al.(2014)]{keller2014}
 Keller, C. J., Honey, C. J., M\'{e}gevand, P., Entz, L., Ulbert, I., \& Mehta, A. D. (2014). Mapping human brain networks with cortico-cortical evoked potentials. \textit{Philosophical Transactions of the Royal Society B}, 369(1653), 20130528. \href{https://doi.org/10.1098/rstb.2013.0528}{doi:10.1098/rstb.2013.0528}.

\bibitem[Kish(2002)]{kish2002}
 Kish, L. B. (2002). End of Moore's law: thermal (noise) death of integration in micro and nano electronics. \textit{Physics Letters A}, 305(3--4), 144--149. \href{https://doi.org/10.1016/S0375-9601(02)01365-8}{doi:10.1016/S0375-9601(02)01365-8}.

\bibitem[Kleiner \& Ludwig(2024)]{kleiner2024}
 Kleiner, J., \& Ludwig, T. (2024). The case for neurons: A no-go theorem for consciousness on a chip. \textit{Neuroscience of Consciousness}, 2024(1), niae037. \href{https://doi.org/10.1093/nc/niae037}{doi:10.1093/nc/niae037}.

\bibitem[Kocher et al.(1999)]{kocher1999}
 Kocher, P., Jaffe, J., \& Jun, B. (1999). Differential power analysis. In \textit{Advances in Cryptology --- CRYPTO '99}, Lecture Notes in Computer Science, vol.~1666, pp.~388--397. Springer. \href{https://doi.org/10.1007/3-540-48405-1_25}{doi:10.1007/3-540-48405-1\_25}.

\bibitem[Kunz \& Pfister(1976)]{kunz1976}
 Kunz, H., \& Pfister, C.-E. (1976). First order phase transition in the plane rotator ferromagnetic model in two dimensions. \textit{Communications in Mathematical Physics}, 46, 245--251. \href{https://doi.org/10.1007/BF01609121}{doi:10.1007/BF01609121}.

\bibitem[Mainen \& Sejnowski(1995)]{mainen1995}
 Mainen, Z. F., \& Sejnowski, T. J. (1995). Reliability of spike timing in neocortical neurons. \textit{Science}, 268(5216), 1503--1506. \href{https://doi.org/10.1126/science.7770778}{doi:10.1126/science.7770778}.

\bibitem[Markram et al.(1997)]{markram1997}
 Markram, H., L\"{u}bke, J., Frotscher, M., Roth, A., \& Sakmann, B. (1997). Physiology and anatomy of synaptic connections between thick tufted pyramidal neurones in the developing rat neocortex. \textit{The Journal of Physiology}, 500(2), 409--440. \href{https://doi.org/10.1113/jphysiol.1997.sp022031}{doi:10.1113/jphysiol.1997.sp022031}.

\bibitem[McFadden(2020)]{mcfadden2020}
 McFadden, J. (2020). Integrating information in the brain's EM field: the cemi field theory of consciousness. \textit{Neuroscience of Consciousness}, 2020(1), niaa016.

\bibitem[Mermin \& Wagner(1966)]{mermin1966}
 Mermin, N. D., \& Wagner, H. (1966). Absence of ferromagnetism or antiferromagnetism in one- or two-dimensional isotropic Heisenberg models. \textit{Physical Review Letters}, 17(22), 1133--1136. \href{https://doi.org/10.1103/PhysRevLett.17.1133}{doi:10.1103/PhysRevLett.17.1133}.

\bibitem[Nyquist(1928)]{nyquist1928}
 Nyquist, H. (1928). Thermal agitation of electric charge in conductors. \textit{Physical Review}, 32(1), 110--113. \href{https://doi.org/10.1103/PhysRev.32.110}{doi:10.1103/PhysRev.32.110}.

\bibitem[Nunez \& Srinivasan(2014)]{nunez2014}
 Nunez, P. L., \& Srinivasan, R. (2014). Neocortical dynamics due to axon propagation delays in cortico-cortical fibers: EEG traveling and standing waves with implications for top-down influences on local networks and white matter disease. \textit{Brain Research}, 1542, 138--166. \href{https://doi.org/10.1016/j.brainres.2013.10.036}{doi:10.1016/j.brainres.2013.10.036}.

\bibitem[Ororbia \& Friston(2023)]{ororbia2023}
 Ororbia, A., \& Friston, K. (2023). Mortal computation: A foundation for biomimetic intelligence. \textit{arXiv:2311.09589}. \href{https://doi.org/10.48550/arXiv.2311.09589}{doi:10.48550/arXiv.2311.09589}.

\bibitem[Pinto et al.(2017)]{pinto2017}
 Pinto, Y., Neville, D. A., Otten, M., Corballis, P. M., Lamme, V. A. F., de Haan, E. H. F., Foschi, N., \& Fabri, M. (2017). Split brain: Divided perception but undivided consciousness. \textit{Brain}, 140(5), 1231--1237. \href{https://doi.org/10.1093/brain/aww358}{doi:10.1093/brain/aww358}.

\bibitem[Prime et al.(2020)]{prime2020}
 Prime, D., Woolfe, M., O'Keefe, S., Rowlands, D., \& Dionisio, S. (2020). Quantifying volume conducted potential using stimulation artefact in cortico-cortical evoked potentials. \textit{Journal of Neuroscience Methods}, 337, 108639. \href{https://doi.org/10.1016/j.jneumeth.2020.108639}{doi:10.1016/j.jneumeth.2020.108639}.

\bibitem[Rabaey et al.(2003)]{rabaey2003}
 Rabaey, J. M., Chandrakasan, A., \& Nikoli\'{c}, B. (2003). \textit{Digital Integrated Circuits: A Design Perspective} (2nd ed.). Prentice Hall.

\bibitem[Rebollo et al.(2021)]{rebollo2021}
 Rebollo, B., Telenczuk, B., Navarro-Guzman, A., Destexhe, A., \& Sanchez-Vives, M. V. (2021). Modulation of intercolumnar synchronization by endogenous electric fields in cerebral cortex. \textit{Science Advances}, 7(10), eabc7772. \href{https://doi.org/10.1126/sciadv.abc7772}{doi:10.1126/sciadv.abc7772}.

\bibitem[Reinhart \& Nguyen(2019)]{reinhart2019}
 Reinhart, R. M. G., \& Nguyen, J. A. (2019). Working memory revived in older adults by synchronizing rhythmic brain circuits. \textit{Nature Neuroscience}, 22(5), 820--827. \href{https://doi.org/10.1038/s41593-019-0371-x}{doi:10.1038/s41593-019-0371-x}.

\bibitem[Rosas et al.(2024)]{rosas2024}
 Rosas, F. E., Geiger, B. C., Luppi, A. I., Seth, A. K., Polani, D., Gastpar, M., \& Mediano, P. A. M. (2024). Software in the natural world: A computational approach to hierarchical emergence. \textit{arXiv:2402.09090}. \href{https://doi.org/10.48550/arXiv.2402.09090}{doi:10.48550/arXiv.2402.09090}.

\bibitem[Sabatini \& Regehr(1996)]{sabatini1996}
 Sabatini, B. L., \& Regehr, W. G. (1996). Timing of neurotransmission at fast synapses in the mammalian brain. \textit{Nature}, 384(6605), 170--172. \href{https://doi.org/10.1038/384170a0}{doi:10.1038/384170a0}.

\bibitem[Schneider et al.(2021)]{schneider2021}
 Schneider, M., Broggini, A. C., Dann, B., Tzanou, A., Uran, C., Sheshadri, S., Scherberger, H., \& Vinck, M. (2021). A mechanism for inter-areal coherence through communication based on connectivity and oscillatory power. \textit{Neuron}, 109(24), 4050--4067. \href{https://doi.org/10.1016/j.neuron.2021.09.037}{doi:10.1016/j.neuron.2021.09.037}.

\bibitem[Seth(2026)]{seth2026}
 Seth, A. K. (2026). Conscious artificial intelligence and biological naturalism. \textit{Behavioral and Brain Sciences}, 49, e315. \href{https://doi.org/10.1017/S0140525X25000032}{doi:10.1017/S0140525X25000032}.

\bibitem[Tononi \& Koch(2015)]{tononi2015}
 Tononi, G., \& Koch, C. (2015). Consciousness: here, there and everywhere? \textit{Philosophical Transactions of the Royal Society B}, 370(1668), 20140167. \href{https://doi.org/10.1098/rstb.2014.0167}{doi:10.1098/rstb.2014.0167}.

\bibitem[Trebaul et al.(2016)]{trebaul2016}
 Trebaul, L., Rudrauf, D., Job, A.-S., M\u{a}l\^{\i}ia, M. D., Popa, I., Barborica, A., Minotti, L., M\^{\i}ndru\c{t}\u{a}, I., Kahane, P., \& David, O. (2016). Stimulation artifact correction method for estimation of early cortico-cortical evoked potentials. \textit{Journal of Neuroscience Methods}, 264, 94--102. \href{https://doi.org/10.1016/j.jneumeth.2016.03.002}{doi:10.1016/j.jneumeth.2016.03.002}.

\end{thebibliography}
